\documentclass[10pt]{article}
\usepackage[utf8]{inputenc}
\usepackage{polski}
\usepackage{float}
\usepackage{amsmath}
\usepackage[margin=0.6in]{geometry}
\usepackage{gensymb}

\begin{document}
{\centering
\title{Sprawozdanie: Analiza porównawcza ludności wdług poziomu wykształcenia w latach 1978 i 2011.}
\date{}
\maketitle}

\section{Wstęp}
Analizowany zbiór danych przedstawia strukturę ludności według wieku oraz poziomu wykształcenia. Dane zostały zestawione dla dwóch okresów oddalonych od siebie o 33 lata (1978--2011), dzięki czemu możliwe jest prześledzenie zmian w poziomie edukacji społeczeństwa. W tabelach wyróżniono grupy wieku oraz poziomy wykształcenia, m.in. wyższe, średnie oraz podstawowe.

\subsection{Źródło danych}

Źródłem danych jest Główny Urząd Statystyczny (GUS), \textit{Ludność według poziomu wykształcenia 1921--2011}.

\begin{flushleft}
	\texttt{https://dane.gov.pl/pl/dataset/1572,ludnosc-wedug-poziomu-wyksztacenia-1921-2011/resource/17481/table}
\end{flushleft}

\subsection{Cel analizy}

Celem analizy jest porównanie powszechności edukacji dawniej i obecnie, czyli określenie, jak zmienił się poziom wykształcenia społeczeństwa na przestrzeni 33 lat. W szczególności analiza ma odpowiedzieć na następujące pytania:

\begin{itemize}
	\item Czy wzrósł udział osób z wykształceniem średnim i wyższym?
	\item Jaki jest przeciętny wiek osoby z wykształceniem podstawowym, średnim oraz wyższym?
	\item Jak zmiany poziomu wykształcenia kształtują się w różnych grupach wiekowych?
	\item Czy edukacja stała się bardziej powszechna i dostępna niż w przeszłości?
\end{itemize}

Przeprowadzenie takiej analizy pozwoli ocenić skalę przemian edukacyjnych w społeczeństwie oraz pokazać, jak bardzo wzrosło znaczenie wykształcenia w ciągu ostatnich dekad.

\section{Analiza rozkładu przy pomocy metryk statystycznych}

\subsection{Średnia arytmetyczna liczności}

\paragraph{Wzór ogólny}
\[
\bar{x} = \frac{x_{1}n_{1} + x_{2}n_{2} + ... + x_{k}n_{k}}{n_{1} + n_{2} + ... + n_{k}} 
\]

\subsection{Średnia ucinana liczności}

\paragraph{Wzór ogólny}
\[
\bar{x} = \frac{x_{i}n_{i} + x_{i+1}n_{i+1} + ... + x_{k-i}n_{k-i}}{n_{i} + n_{i+1} + ... + n_{k-i}} 
\]


\subsection{Mediana}

\paragraph{Wzór ogólny}
\[
x_{med} = 
    \begin{cases}
        x_{(\frac{k+1}{2})} \text{ dla k nieparzystych}\\
        \frac{1}{2} (x_{(\frac{k}{2})} + x_{(\frac{k+1}{2})}) \text{ dla k parzystych}
    \end{cases}
\]

\subsection{Kwartyle}
\paragraph{Wzory ogólne}
\paragraph{Q1}
\[
Q_{1} =
    \begin{cases}
        x_{(\frac{k+1}{4})} & \text{gdy } \frac{k+1}{4} \in N \\
        (1-\alpha)\,x_{(\lfloor\frac{k+1}{4}\rfloor)} + \alpha\,x_{(\lceil\frac{k+1}{4}\rceil)}
        & \text{gdy } \alpha = \frac{k+1}{4} - \lfloor\frac{k+1}{4}\rfloor
    \end{cases}
\]

\paragraph{Q2}
\[
Q_{2} = x_{med} =
    \begin{cases}
        x_{(\frac{k+1}{2})} & \text{dla } k \text{ nieparzystego} \\
        \frac{1}{2} \left( x_{(\frac{k}{2})} + x_{(\frac{k}{2}+1)} \right)
        & \text{dla } k \text{ parzystego}
    \end{cases}
\]

\paragraph{Q3}
\[
Q_{3} =
    \begin{cases}
        x_{(\frac{3(k+1)}{4})} & \text{gdy } \frac{3(k+1)}{4} \in N \\
        (1-\beta)\,x_{(\lfloor\frac{3(k+1)}{4}\rfloor)} + \beta\,x_{(\lceil\frac{3(k+1)}{4}\rceil)}
        & \text{gdy } \beta = \frac{3(k+1)}{4} - \lfloor\frac{3(k+1)}{4}\rfloor
    \end{cases}
\]

\subsection{Rozstęp międzykwartylowy}
\paragraph{Wzór ogólny}
\[
IQR = Q_{3} - Q_{1}
\]

\subsection{Rozstęp}
\paragraph{Wzór ogólny}
\[
R = x_{(k)} - x_{(1)}
\]

\subsection{Wariancja i odchylenie standardowe}
\paragraph{Wzory ogólne}
\paragraph{Wariancja}
\[
\sigma^{2} = \frac{1}{k} \sum_{i=1}^{k} (x_{i} - \bar{x})^{2}
\qquad
s^{2} = \frac{1}{k-1} \sum_{i=1}^{k} (x_{i} - \bar{x})^{2}
\]

\paragraph{Odchylenie standardowe}
\[
\sigma = \sqrt{\sigma^{2}}
\qquad
s = \sqrt{s^{2}}
\]

\subsection{Współczynnik zmienniczości}
\paragraph{Wzory ogólne}
\[
CV = \frac{\sigma}{|\mu|} \cdot 100\%
\qquad
CV = \frac{s}{|\bar{x}|} \cdot 100\%
\]

\subsection{Skośność i kurtoza}

\paragraph{Wzory ogólne}

\paragraph{Współczynnik skośności}
\[
\gamma_{1} = \frac{\frac{1}{k} \sum_{i=1}^{k} (x_{i} - \bar{x})^{3}}
                         {\left(\frac{1}{k} \sum_{i=1}^{k} (x_{i} - \bar{x})^{2}\right)^{3/2}}
\]

\paragraph{Kurtoza}
\[
\gamma_{2} = \frac{\frac{1}{k} \sum_{i=1}^{k} (x_{i} - \bar{x})^{4}}
                         {\left(\frac{1}{k} \sum_{i=1}^{k} (x_{i} - \bar{x})^{2}\right)^{2}}-3
\]



\section{Prezentacja rozkładów na wykresach}

\section{Podsumowanie}

Przeprowadzona analiza porównawcza danych z lat 1978 oraz 2011 wykazała istotne zmiany w strukturze wieku osób należących do poszczególnych poziomów wykształcenia. Zaobserwowane różnice dotyczą zarówno przeciętnego wieku badanych grup, jak i stopnia zróżnicowania danych oraz kształtu analizowanych rozkładów.

\subsection{Osoby z wykształceniem wyższym}

W 1978 roku średni wiek osób z wykształceniem wyższym wynosił 37{,}30 roku, natomiast w 2011 roku wzrósł do 40{,}83 roku, co oznacza wzrost o około 3{,}5 roku. Mediana wieku zwiększyła się z 34 do 36 lat, co wskazuje na przesunięcie całego rozkładu ku starszym grupom wieku.

Wartość współczynnika skośności wzrosła z 0{,}49 w 1978 roku do 0{,}90 w 2011 roku, co oznacza silniejszą asymetrię prawostronną rozkładu. W praktyce wskazuje to na koncentrację obserwacji wokół młodszych grup wieku przy jednoczesnym występowaniu większej liczby obserwacji w starszych grupach wiekowych. Współczynnik zmienności wzrósł z 33{,}01\% do 41{,}43\%, co świadczy o zwiększeniu zróżnicowania wieku w tej grupie.

wskazuje to na rosnącą liczebność osób posiadających wyższe wykształcenie oraz zwiększenie dostępności edukacji wyższej.

\subsection{Osoby z wykształceniem średnim}

W przypadku osób z wykształceniem średnim zaobserwowano jeden z największych wzrostów przeciętnego wieku. Średnia wieku zwiększyła się z 34{,}12 roku w 1978 roku do 43{,}20 roku w 2011 roku, czyli o ponad 9 lat. Mediana wzrosła z 36 do 42 lat, natomiast trzeci kwartyl zwiększył się z 44 do 56 lat, co wskazuje na istotne przesunięcie struktury wieku ku starszym grupom.

W 1978 roku rozkład tej grupy charakteryzował się wyraźną asymetrią prawostronną, czego potwierdzeniem jest współczynnik skośności równy 1{,}52. W 2011 roku wartość ta spadła do 0{,}37, co oznacza znaczne zmniejszenie asymetrii oraz bardziej równomierny rozkład wieku. Dodatkowo wartość kurtozy zmieniła się z 3{,}55 do $-0{,}87$, co wskazuje na przejście od rozkładu bardziej skoncentrowanego wokół średniej do rozkładu bardziej spłaszczonego.

Otrzymane wyniki sugerują, że grupa osób z wykształceniem średnim uległa znacznemu starzeniu, a jednocześnie rozkład wieku stał się bardziej stabilny i zrównoważony.

\subsection{Osoby z wykształceniem podstawowym}

Największe zmiany zaobserwowano w grupie osób z wykształceniem podstawowym. Średni wiek wzrósł z 44{,}99 roku w 1978 roku do 57{,}71 roku w 2011 roku, czyli o prawie 13 lat. Mediana wieku zwiększyła się z 47 do 62 lat, co świadczy o bardzo silnym przesunięciu rozkładu ku starszym grupom wiekowym.

Istotna zmiana nastąpiła również w kształcie rozkładu. W 1978 roku współczynnik skośności wynosił 0{,}66, co wskazywało na asymetrię prawostronną. W 2011 roku przyjął on wartość ujemną równą $-0{,}68$, co oznacza odwrócenie kierunku asymetrii i koncentrację obserwacji w starszych grupach wieku. Jednocześnie współczynnik zmienności zmniejszył się z 37{,}29\% do 34{,}44\%, co wskazuje na większe skupienie obserwacji wokół średniej.

Uzyskane wyniki pokazują, że w 2011 roku wykształcenie podstawowe stało się cechą charakterystyczną przede wszystkim dla starszych pokoleń, podczas gdy w młodszych grupach wieku jego udział wyraźnie zmalał.

\subsection{wnioski}

Przeprowadzona analiza porównawcza wykazała wyraźne przemiany struktury edukacyjnej społeczeństwa pomiędzy rokiem 1978 a 2011. We wszystkich analizowanych grupach odnotowano wzrost przeciętnego wieku, jednak skala tych zmian była zróżnicowana. Największe przesunięcie ku starszym grupom wieku nastąpiło wśród osób z wykształceniem podstawowym oraz średnim, natomiast grupa osób z wykształceniem wyższym zachowała relatywnie młodszą strukturę wieku.

Najbardziej znaczącą zmianą jakościową było przekształcenie grupy osób z wykształceniem podstawowym w populację silnie skoncentrowaną w starszych grupach wiekowych. Jednocześnie zwiększenie udziału osób z wykształceniem średnim i wyższym w młodszych grupach wieku wskazuje na wzrost dostępności edukacji oraz systematyczne podnoszenie poziomu wykształcenia społeczeństwa.

Wyniki analizy statystycznej potwierdzają, że w analizowanym okresie nastąpiła istotna transformacja społeczna i edukacyjna, prowadząca do stopniowego ograniczania udziału osób z najniższym poziomem wykształcenia oraz wzrostu znaczenia wykształcenia średniego i wyższego w strukturze społecznej.


\end{document}